% ============================================================
% cita-ieee-html.tex — IEEE EN PANTALLA (SC-57)
% ------------------------------------------------------------
% LO CARGA SOLO EL DOCUMENTO MÍNIMO DE LA PASADA DE tex4ht,
% DESPUÉS DE cita-ieee-config.tex. CON ESTE ARCHIVO, HTML Y EPUB
% TOMAN EL TEXTO DE LAS CITAS Y DE LAS REFERENCIAS DE biblatex-ieee
% (m_GenerarSalidasLibro.PrepararCitasPantalla).
%
% IEEE NO CAMBIA NADA EN PANTALLA: cita-ieee-config.tex YA CARGA
% backref=false, Y EL ESTILO NO USA VERSALITAS. LA ETIQUETA [n] DE
% CADA ENTRADA LLEGA EN EL <dt> DE citas.html Y LA PONEN LAS HOJAS
% (cita-fragmentos.xsl). LA NUMERACIÓN VUELVE A 1 EN CADA PIEZA: CADA
% UNA ES UNA refsection, COMO EN EL PDF. MEDIDO CON EL LIBRO DE DUEK
% (132 CITAS, 113 REFERENCIAS): EL MISMO TEXTO QUE EL PDF.
%
% LAS NOTAS: biblatex NO LAS DETECTA BAJO tex4ht («Patching
% footnotes failed», SC-57). \iffootnote SOLO LO USA \smartcite DE
% numeric-comp.cbx, QUE LA APLICACIÓN NO EMITE (biblatex-ieee v1.3f,
% RF-14).
% ============================================================

% ------------------------------------------------------------
% \gbImprimirBibliografia — LA EMITE docbook-to-latex-citas.xsl
% ------------------------------------------------------------
% SIN TÍTULO: EL DE LA BIBLIOGRAFÍA LO PONE LA HOJA DE HTML O EPUB
% ------------------------------------------------------------
\newcommand*{\gbImprimirBibliografia}{%
  \printbibliography[heading=none]}
